% =====================================================================
%  Harald Høffding, Gensvar til Dr. Georg Brandes
%  In: Tilskueren. Maanedsskrift for Literatur, Samfundsspørgsmaal og almenfattelige videnskabelige Skildringer,
%  udgivet af N. Neergaard.  Syvende Aargang.  Kjøbenhavn: P. G. Philipsens Forlag; Thieles Bogtrykkeri.  Februar--Marts 1890, pp. 125-153.
%  Antiqua.
%
%  TRANSCRIPTION -- GENERATED; DO NOT EDIT BY HAND.  Høffding's four pieces in Tilskueren 1886-1890 (Om Religionsfrihed;
%  the three contributions to the Nietzsche debate with Georg Brandes) are built by ONE pipeline
%  (ebook-method/tilskueren-1886-90/) into a combined working file, which split.py cuts into four transcriptions.
%  Rebuild from the repo root:   sh ebook-method/tilskueren-1886-90/rebuild.sh
%
%  STATUS: COMPLETE (2026-10-15): the volume's title page and the whole text, pp. 125-153, with its notes.
%
%  EDITION.  The only printing in Tilskueren.  Text of record: Google's scan of the Harvard copy, Internet Archive
%  tilskuerenmaaned1890nngoog, 600 ppi grey JP2 leaves with the Archive's ABBYY OCR; extract ~/bibliotek/Høffding, Harald/1890-tilskueren-7-pp125-153-343-353.pdf (SCANS.tsv, sha256),
%  made by ia_extract.py (title leaf, Indhold leaves, the article).  Page range confirmed from the volume's Indhold and the
%  signature.  Høffding's answer to Brandes's »Det store Menneske, Kulturens Kilde« (Tilskueren, January 1890).  Translated by William Banks, The Great Debate (2024); not translated here.  The print spells Nietzsche so here (the 1889 piece has Nietsche).
%
%  ACCURACY (2026-10-15; one collation for all four texts, 76 pp., ~29,000 words).
%  (a) The DRAFT is the Archive's ABBYY OCR (the layer of the working PDF), turned into LaTeX by draft.py; it was collated
%      word by word with an independent tesseract reading (dan, tessdata_best) of the same images (collate/check.py,
%      witness 'ocr'): 650 disagreements.  check.py's own rules settled 44.
%  (b) A THIRD reading voted: Project Runeberg's OCR of the same Google images (never proofread there), by context
%      (rules.py umn-*, 372 decisions), guarded so that a vote never makes a word the corpus does not attest; other
%      rules (the layer's glyph faults against an attested word, accents tesseract adds, the ø rule, the house thin space)
%      settled 70.  The vote was AUDITED at the image: 40 rule decisions sampled blind and 30 more suspects read at zooms:
%      7 Runeberg-for-the-draft votes were wrong (ABBYY and Runeberg sharing a fault: bg/og, kim/kun, bar/har,
%      Anstrengelser/Anstrængelser, Vist/vist, én/en, h Vad/hvad), all mended; the ~30 letter-changing votes for the draft
%      not looked at are tesseract noise against a real word.
%  (c) The rest were read at the image BLIND (round b1: 303 crops on 31 sheets, six readers, 62/62 planted controls right):
%      203 open checks, 23 places where Runeberg disputed both OCRs, 54 places where the henrikpontoppidan.dk e-text
%      differs (a FINDING AID only: someone else's transcription, never a vote -- it normalises the print's spelling, e.g.
%      Nietzsche for the 1889 print's Nietsche), and the audit; ~150 decided in the main session from the readers' lines,
%      ~25 at 300-dpi zooms.  OPEN 0.  An unattested-word sweep of the collated text (353 words in no other transcription)
%      was read through; its doubtful words were looked at in zooms (one misprint found: forkellige).
%  (d) STRUCTURE read at the image (reader K): the title pages, the article heads, the section numerals (two the draft
%      had lost), the verse, the divider; every FOOTNOTE read whole (reader N: 11 notes; compared word by word with the
%      collated notes: 0 word differences, the reader right on the closing quotes).
%  (e) LETTERSPACING (the print's emphasis, set \emph{} here) is MEASURED at the image, not read: two blind line rounds
%      (l1, 256 lines; l2, the 134 doubtful ones again) agreed on only 40 of 134 lines, so spacemeasure.py decides by the
%      white between the letters of each word in tesseract's boxes (bimodal: 1-4 px roman, 6-8 px spaced; threshold 5 px;
%      isolated words must have every gap wide).  Checked at zooms on 16 words where the measure and the e-text's <em>
%      differ: the measure right on 11; three of its false positives (heavy isolated words: Ævne, frem, Side) led to the
%      isolation rule; three misses mended by hand after the page check (SPACEFIX in spaceapply.py).
%  (f) SAMPLED PAGE CHECK (pagecheck.py, seed 20261015; two blind readers): 6 whole pages, one or two from each piece
%      (Religionsfrihed p. 451; Demokratisk Radikalisme pp. 852, 868; Gensvar pp. 132, 140; Epilog p. 343), 2202 words:
%      1 WORD error (Talemaader for the print's Talemaaler: a correction wrongly taken from the e-text finding aid),
%      1 PUNCT (a speck read as '), both mended; LETTERSPACING 15 differences: at 300-dpi zooms 5 transcription errors
%      (tre, Kopier twice, Genmæle, fortsatte; mended), 4 where the transcription is right, 6 not decidable at 300 dpi
%      (the ends of runs in wide-set lines).  Word error rate after mending 0 in 2202 (95% upper bound ~1.4 per 1000);
%      letterspacing remains the least certain feature: ~2-3 per 1000 words.
%  PRINTED AS IS (misprints kept, marked at the line):
%    nothing recorded.
%  CONVENTIONS.  Page turns: % --- p. N --- and \opage{N} (the printed page); a word divided over a page turn is split with
%  %.  Footnotes as \footnote at their marks, numbered per page as printed (1)).  \emph{} = letterspacing; signatures in
%  small capitals as printed.  The running heads and the sheet signatures are not transcribed.  The ø rule (CLAUDE.md).
%  Quotation marks as printed: »...« and „...“ both occur, and the print sometimes doubles them (»»Zunächst).
% =====================================================================
\documentclass[12pt,a4paper,openany]{book}
\usepackage[utf8]{inputenc}
\usepackage[T1]{fontenc}
\usepackage{libertinus}
\usepackage{libertinust1math}
\usepackage[danish]{babel}
\usepackage[protrusion=true,expansion=true]{microtype}
\usepackage{setspace}
\usepackage[margin=1.3in]{geometry}
\usepackage{fancyhdr}
\usepackage{hyperref}
\hypersetup{pdftitle={Gensvar til Dr. Georg Brandes}, pdfauthor={Harald Høffding}, hidelinks}
\pagestyle{fancy}
\fancyhf{}
\fancyhead[LE,RO]{\thepage}
\fancyhead[C]{\textit{H. Høffding: Gensvar til Dr. Georg Brandes (1890)}}
\fancypagestyle{plain}{\fancyhf{}\fancyhead[LE,RO]{\thepage}}
\setstretch{1.3}
\usepackage{marginnote}
\usepackage{graphicx}
\DeclareUnicodeCharacter{0254}{\reflectbox{c}}   % ɔ (id est)
\renewcommand{\thefootnote}{\arabic{footnote})}   % the print numbers each page's notes 1), 2)
\newcommand{\opage}[1]{\marginnote{\footnotesize\textit{[#1]}}}
\newcommand{\sectionnum}[1]{\par\medskip\begin{center}#1\end{center}\smallskip}
\newcommand{\asterism}{\par\begin{center}*\quad*\\[-0.4em]*\end{center}}
\newcommand{\signature}[1]{\par\smallskip\hfill #1\hspace*{2em}\par}
\pdfstringdefDisableCommands{\renewcommand{\footnote}[2][]{}\renewcommand{\opage}[1]{}}

\begin{document}

\begin{titlepage}
\begin{center}
\vspace*{2em}
{\LARGE\bfseries \emph{TILSKUEREN.}}\\[1em]
{\large \emph{MAANEDSSKRIFT}}\\[0.5em]
{\small FOR}\\[0.5em]
LITERATUR, SAMFUNDSSPØRGSMAAL OG ALMENFATTELIGE\\
VIDENSKABELIGE SKILDRINGER.\\[1.5em]
{\small UDGIVET AF}\\[0.5em]
{\bfseries N. NEERGAARD.}\\[1.5em]
\emph{SYVENDE AARGANG.}\\
\vfill
KJØBENHAVN 1890.\\[0.4em]
{\bfseries P. G. PHILIPSENS FORLAG.}\\[0.3em]
{\small THIELES BOGTRYKKERI.}
\end{center}
\end{titlepage}
\clearpage

% --- p. 125 ---
\setcounter{footnote}{0}%
\opage{125}
\begin{center}{\large\bfseries Gensvar til Dr. Georg Brandes.}\end{center}\medskip

\sectionnum{I.}

Der er et forholdsvis stort Parti af Dr. Georg Brandes
Artikel i »Tilskuerens« Januarhæfte, som jeg kun ugærne svarer
paa. Det er den almindelige Karakteristik af mig, han giver som
Indledning. Jeg vil helst gøre, som naar jeg bliver fotograferet,
nemlig overlade andre at dømme om Ligheden. Dr. Brandes
beraaber sig paa, at jeg selv har tillagt ham en Ævne til »psykologisk
Karakteristik«. Han har ikke lagt Mærke til mine Ord; jeg sagde
med velberaad Hu »\emph{stemningsfuld} Karakteristik«. Den Stemning,
der præger hans Karakteristiker, stammer altfor ofte fra hans rent
personlige, forud tilfældigt vakte Antipati eller Sympati overfor sin «
Genstand. Jeg anerkender i denne Henseende hans Karakteristik
af mig som \emph{meget} stemningsfuld. Men med Hensyn til vore Stridspunkter
er det ligegyldigt, om han har Ret i denne Karakteristik
eller ikke. Min Ringhed kunde være endnu ringere end den er,
og dog kunde jeg have Ret i mit Angreb paa den »aristokratiske
Radikalisme«. Blind Høne kan ogsaa finde et Korn. Men det er
et Træk, der ikke er ukendt hos Dr. G. Brandes som Polemiker,  
at aflede Opmærksomheden fra Hovedsagen ved uvedkommende
Bemærkninger, især af personlig Art. --- Der er alligevel enkelte »
Punkter i denne Del af Artiklen i Januarhæftet, som jeg ikke
kan forbigaa.

Der er vel næppe nogen Klasse af Mennesker, som Dr. G.
Brandes mener at staa fjærnere, end Teologerne. Og dog har
han i sin Maade at polemisere paa ikke lidet fælles med dem.
Han véd i Regelen lige saa godt som saa mange Teologer særdeles
god Besked med sine Modstanderes Motiver, og som saa mange
% --- p. 126 ---
\setcounter{footnote}{0}%
\opage{126}Teologer er han i Regelen ogsaa sikker paa, at det er mindre
ædle Motiver, der leder dem. Det Angreb, jeg har rettet mod
den aristokratiske Radikalisme, udleder han af, at jeg har »en
Tilbøjelighed til at tage dygtig Afstand fra Personligheder, hvis
Lod i Livet er bleven den, altid at maatte føre den upopulære
Sag«. Den Insinuation, som heri ligger, afviser jeg som usømmelig
og som usand. Jeg har selv gjort Rede for, hvad der bestemte
mig til at skrive min Artikel »Demokratisk Radikalisme«, og jeg
kan forlange at blive troet paa mit Ord. Uagtet jeg straks ved
Læsningen af Artiklen »Aristokratisk Radikalisme« følte mig frastødt
af det selvmodsigende og usunde i dens Tankegang., betænkte
jeg mig dog længe paa at træde op med en offentlig Kritik, fordi
jeg ikke vilde angribe en fortjænstfuld Mand, der har saa mange
Angribere, med hvilke jeg ikke vilde sættes i Klasse, og af hvilke
jeg ikke vilde tages til Indtægt. Grunden til, at jeg alligevel gjorde
det, var ikke den, at Dr. G. Brandes er \emph{upopulær}, men netop
den, at han er populær. Det er noget, som han under sine
hyppige Udeladelser om sin Forfølgelse og sin Upopularitet er tilbøjelig
til at overse, at han har en stor Tilhængerkreds\footnote[1]{Hvilken Slags Mennesker denne Kreds bestaar af, eller rettere Dr Brandes maa ønske, at den skal bestaa af, kan maaske være noget vanskeligt at sige. Da han med en vis Haan overlader mig „Tidsalderens \emph{dannede} Kredse“, kan det ikke være dannede, hin Kreds bestaar af. Da det vel heller ikke kan være udannede, bliver der ikke andre end de \emph{overdannede} tilbage. Mindre kan vel heller ikke gøre det.}, hvad
der jo heller ikke kan undre, da han har den lykkelige Gave at
kunne tale med Styrke og Begejstring om hvad han i Øjeblikket
er optagen af, og derfor har saa let ved at skaffe Ørenlyd for
hvad han vil have frem. Jeg mærkede da ogsaa, at hin Artikel
havde øvet Indflydelse paa ikke faa. Derfor bestemte jeg mig til
at skrive, medens jeg først havde tænkt at indskrænke mig til at
behandle Spørgsmaalet i en Kreds af Studerende, der samles hos
mig. Jeg vidste meget vel, at Dr. Brandes selv og hans Ildtilbedere
vilde tage mig det ilde op, og at de nok skulde regne ud, at jeg
havde gjort det af et eller andet mindre smukt Motiv. Men Spørgsmaalet
var, om Dr. Brandes, fordi han har været Genstand for
saa mange uretfærdige og fanatiske Angreb, derfor slet ikke skulde
kunne kritiseres, selv hvor han gaar ind paa andre Omraader end
hans eget. At' sætte ham udenfor enhver Kritik vilde i mine Øjne
% --- p. 127 ---
\setcounter{footnote}{0}%
\opage{127}være at gøre ham en meget daarlig Tjæneste. Er der noget, der
har været til hans Skade, saa er det det, at han ikke i Tide har
været underkastet tilstrækkelig skarp \emph{og tillige} tilstrækkelig
upartisk Kritik.

At Dr. Brandes ikke har Ret til at tillægge mig en Tilbøjelighed
til at tage Afstand fra ham som upopulær, kan jeg
godtgøre ved et bestemt Faktum. Jeg har aldrig omtalt Dr. Brandes
hverken i mine Skrifter eller i mine Forelæsninger (uden enkelte
Steder, hvor jeg benytter nogle af hans literaturhistoriske Skrifter).
Kun én eneste Gang har jeg omtalt ham offentlig. For nogle
faa Aar siden kom jeg til at fungere som officiel Opponent ved
en literaturhistorisk Disputats. Jeg begyndte da mit Foredrag
med at beklage, at det udmærkede Talent, som vort Land ejede
paa den æstetiske Kritiks Omraade, ikke ved en Lejlighed som
denne kunde føre Ordet paa Universitetets Vegne. Dette kan dog
vel ikke kaldes at tage Afstand?

Jeg skal ikke nægte, at der siden den Tid har været adskilligt
i Dr. Brandes' offentlige Udtalelser og Optræden, der saare
lidet har tiltalt mig og har gjort, at jeg \emph{føler mig} i Afstand fra
ham. Naar jeg ikke er traadt op imod ham, er det dels af den
ovenanførte Grund, dels fordi jeg i det hele helst vil høre til de
stille i Landet og passe mine egne Sager.

I øvrigt forekommer det mig underligt, at Dr. Brandes i
samme Øjeblik som han tillægger mig en Tilbøjelighed til at give
»ædle Roller«, tilskriver sig selv »den Lod at føre den upopulære
Sag«. Det er dog en ret ædel Rolle, han derved fremtræder i,
medens han tillægger mig den mindre ædle at jage efter Popularitet.
---

Dr. Brandes minder om, at han allerede for 24 Aar siden
havde et Sammenstød med mig, og fremhæver, at jeg har forladt
det Standpunkt, hvorpaa jeg den Gang stod. Heller ikke dette
kommer den Sag ved, Striden gælder. Jeg kan have haft Uret
for 24 Aar siden og dog have Ret nu. Men da der er mindet
om min tidligere Anskuelse, vil jeg bemærke, at jeg selv stadig
har en Følelse af, at der er en indre Sammenhæng mellem min
nuværende Anskuelse og den, jeg i min første Ungdom sluttede
mig til. Jeg har altid senere følt det som en Styrke, at jeg gav
mig Tid til saa alvorligt og grundigt, som det var mig muligt, at
gennemleve og gennemtænke det teologiske Stadium; med hvilket
de fleste begynde deres aandelige Udvikling, hvor denne saa end
% --- p. 128 ---
\setcounter{footnote}{0}%
\opage{128}senere fører dem hen. Jeg klamrede mig saa fast til det, at jeg
for ikke at komme til at opgive det, greb en, som jeg senere
indsaa, umulig Udvej, Distinktionen mellem Tro og Viden, saaledes
som den lærtes af S. Kierkegaard og R. Nielsen. Gennem
den Tids indre Erfaringer og Kampe mener jeg at have vundet
en Sans for Betydningen af personligt, inderligt Liv og en Forstaaelse
af det religiøse Liv, som ingen Kritik og ingen videnskabelig
Analyse senere har svækket, saa meget end min egen
Tro har ændret sig. Dogmernes ubetingede Uholdbarhed blev for
mig kun den ene Side af Sagen. Den Overbevisning har dannet
sig hos mig, at det religiøse Problem først stiller sig i sin fulde
Dybde, naar man gaar nærmere ind paa Betydningen af det
menneskelige Følelsesliv, som har affødt Dogmerne, og som har
haft sin Næring og sit Tilhold i dem. Den Forandring, Følelseslivet
vil undergaa, naar dette dogmatiske Tilhold helt falder bort,
kan man endnu næppe danne sig nogen Forestilling om. Baade
fra konservativ og fra radikal Side er man dog, tror jeg, tilbøjelig
til at overdrive Forandringerne, --- hist, fordi man er vant til at
knytte alle Livsværdier til de dogmatiske Forestillinger, --- her,
fordi den polemiske Holdning overfor Dogmerne let udvider sig til
ogsaa at gælde Følelser, der en Stund fik Luft i Dogmerne, men
i sig selv ville \emph{kunne} faa Luft paa anden Vis. Særligt er der
overordentlig meget af Kristendommens etiske Indhold, som kan
beholde sin Værdi, selv om de dogmatiske Støtter, det hidtil
hvilede paa, faldt bort. I den aandelige Udvikling virker Motivforskydningens
Lov, i Følge hvilken Følelser og Institutioner, som
Menneskeslægten er kommen til og har levet sig ind i under visse
Motivers Indflydelse (her: under Indflydelse af Troen paa det
overnaturlige), kunne bevare deres Liv og deres Betydning under
Indflydelse af helt andre Motiver end de oprindelige (her: den
rent humane Sympati og Indsigten i Betingelserne for Slægtens
sunde, harmoniske Udvikling). Denne Lov er paa det etiske og
sociale Omraade, hvad Loven om Energiens Bestaaen er paa det
fysiske Omraade. Naturligvis gælder denne Parallel kun med den
Forskel, som Forskellen mellem en etisk-social og en fysisk Lov
medfører. Men saa meget staar fast: Menneskene have i deres
religiøse Forestillinger udtrykt nogle af deres allervæsentligste
Livserfaringer, som ikke \emph{behøve} at gaa tabt, og som ikke \emph{bør}
gaa tabt for Slægtens fortsatte Liv, selv om deres hidtilværende
Iklædning falder bort og maa falde bort. --- Jeg lægger overalt
% --- p. 129 ---
\setcounter{footnote}{0}%
\opage{129}stor Vægt paa Kontinuiteten, som jo er en af den nyere Tænknings
vigtigste Principer, og hvis Betydning paa det aandelige
Omraade ikke er mindre end i den fysiske Natur, selv om den er
vanskeligere at paavise og fastholde hist end her

Jeg er sikkert ved denne Overbevisning ikke kommen de
ortodokse nærmere. Under en Diskussion om religiøse Spørgsmaal,
der for nogle Aar siden holdtes i Studenterforeningen, mindede
jeg paa given Anledning spøgende om, at jeg selv havde været
Teolog og tilføjede, at hvis det endnu gjaldt, hvad man sagde i
min Tid: pectus facit theologum (det er Hjærtet, der gør Teologen),
vilde jeg for saa vidt gærne være det endnu. Det blev mig senere
fortalt, at en ung Teolog blandt de tilstedeværende sagde til
sin Kammerat: »Det skulde han dog helst have ladet være at
minde om!«

Den Overbevisning, jeg saaledes mener at have vundet, gør,
at jeg ikke med samme Glæde og Frejdighed som mange andre
ser paa Nedrivelsens Værk. Jeg véd meget vel, at jeg ikke er
nogen revolutionær Natur, og jeg \emph{vil} heller ikke være det. (Det
er aldeles uberettiget, naar Dr. Brandes faar ud af en af mine
Ytringer, at jeg skulde have betegnet mig selv som hørende til
de revolutionære Aander). Jeg er Evolutionist. Først da jeg
lærte den moderne Udviklingslære at kende, følte jeg fast Grund
under Fødderne. Lige saa vis, som jeg var bleven paa, at den
gamle Tro ikke kunde hævdes, lige saa vis blev jeg nu paa, at
den nye Tro maatte udvikle sig af Spirer, der allerede var til
Stede, og at disse Spirer ikke kunde udfoldes uden gennem lige
saa alvorlige og dybtgaaende Livserfaringer som de, gennem
hvilke den gamle Tro var gaaet over i Slægtens Kød og Blod.
Jeg maa tilstaa, at jeg ikke med særlig Tillid og Tryghed har
fulgt den fritænkeriske Bevægelse her hjemme i de senere Aartier.
Den har, forekommer det mig, haft for overfladisk en Karakter,
har opereret altfor meget med færdige Resultater, med Meninger,
der laa i Luften, og den har hos den enkelte ikke altid baaret
Vidne om at være Frugten af personlig Erfaring og selvvirksom
Tænkning. Ved Siden af negativ Kritik har en æstetiserende
Tendens staaet i første Linje, og begge have ofte haft en vis
Blaserthed med sig som naturlig Følgesvend.

Heldigvis kender jeg mange dygtige unge Mennesker, paa
hvem denne Skildring ikke passer. Men den passer paa ikke faa,
og jeg ser her en alvorlig Fare for den aandelige Udviklings
% --- p. 130 ---
\setcounter{footnote}{0}%
\opage{130}Sundhed. Jeg vil ingenlunde tilsløre \emph{den} store Modsætning mellem
den rent menneskelige Livsopfattelse og den positivt religiøse.
Kun tror jeg, at mange have godt af at mindes Lessings Ord:
»Ikke alle ere frie, som driver Spot med deres Lænker!« ---

Jeg har altsaa intet imod at blive mindet om, at jeg for
24 Aar siden stod paa et andet Standpunkt end nu. Dr. Brandes
synes rigtigt nok deraf at drage den Slutning, at da jeg har forladt
hint tidligere Standpunkt, vil jeg ogsaa kunne forlade mit
nuværende. Jeg indrømmer Muligheden, for saa vidt jeg haaber
at lære saa længe jeg lever. Det almindelige Standpunkt, paa
hvilket jeg nu staar, og fra hvilket jeg har forfattet mine vigtigste
Skrifter, naaede jeg til for omtrent 16 Aar siden, da jeg havde
fuldendt min Orientering paa det filosofiske Omraade. Skal jeg
forlade det, haaber jeg, at det først vil ske, efterat jeg har uddybet
det adskilligt mere end jeg endnu har formaaet.*

Hvad Dr. Brandes angaar, er han nu naaet til et meget
ophøjet Standpunkt. Han er (ganske vist i Skøderne paa Nietzsche)
naaet til Overmennesket. Jeg svimler ved Tanken om, hvor
højt han vil naa, naar han gør sin næste Svingning. For mit
eget Vedkommende har jeg foreløbig nok at gøre med det menneskelige.
Kun gennem Uddybelse og Udvidelse af det gaar den Vej
fremad, som jeg kan faa Øje paa.

\sectionnum{II.}

Der er intet, Filosofien og særligt Etiken trænger mere til
i vor Tid end en kritisk og alsidig Afvejen af Principernes Gyldighed
og Rækkevidde, en Prøvelse af deres Anvendelighed paa
saa mange forskellige Forhold som muligt. Den, hvem sligt ikke
interesserer, viser, at han mangler filosofisk Interesse, men derfor
kan han, som jeg har anerkendt for Nietzsches Vedkommende,
meget godt være en fremragende Forfatter, naar man betragter
ham fra andre Synspunkter. Af Aandrigheder, der blænde et
Øjeblik, men briste som Sæbebobler ved nærmere Berøring, har
Filosofien i en foregaaende Periode i Tyskland lidt mere end nok.
De Nietzscheske Teorier, som Dr. G. Brandes paa en meget ukritisk
Maade har gjort sig til Talsmand for, falde ganske ind under
denne Rubrik.

Dr. Brandes siger, at mine Indvendinger mod Nietzsches
Lære om de store Mænd »tage sig ud, som hørte man Moralisten
% --- p. 131 ---
\setcounter{footnote}{0}%
\opage{131}paa et borgerligt Skuespil, opfyldt af det til Borgerheld sigtende,
rette sin Indsigelse mod en patetisk eller spottende Replik af
en Prometheus, Hamlet eller Faust«. Denne Bemærkning har
interesseret mig, fordi jeg selv oftere i Samtale om Nietzsche har
gjort en lignende Betragtning gældende. Min Mening var: Nietzsches
Tankegang kunde være af stor Interesse, og det kunde aldrig falde
mig ind at rette nogen Kritik imod den, hvis den var kommen
frem i Form af Replik eller Monolog i et Drama. Jeg vilde ikke
en Gang have kritiseret ham, hvis han ikke i Artiklen »Aristokratisk
Radikalisme« var bleven præsenteret som en stor Filosof. Hvis
det, han giver, skal være Filosofi, er det meningsløst. Filosofien
er nu én Gang mere prosaisk end Digtekunst og Retorik, og den
har maaske næppe endnu lært at være saa prosaisk, som den
bør være.

Jeg tror at kunne sige, at jeg ikke overvurderer Etikens
Videnskabelighed. Der vil ikke kunne paavises nogen Fremstilling
af Etiken, som er saa lidet dogmatisk, som min. Jeg har ladet
mig det være særligt paaliggende at vise Grænserne for Begrundelse
paa dette Omraade, baade hvad Opstillingen af Principerne, og
hvad deres Anvendelse angaar. Og jeg har i Fortalen udtrykkelig
udtalt Muligheden af, at jeg, trods min Bestræbelse, ikke overalt
har naaet den fulde Begrundelse af mine Paastande. Jeg har altsaa
selv advaret mine Læsere. Der er ingen, som har forlangt af
Dr. Brandes, at han skulde lade sig imponere.

En underlig Udtalelse er denne: »Da jeg endnu har en lille
Tvivl, om Moralen virkeligt er en Videnskab, saa betragter jeg de
Principer, som nogle af de bedste Hoveder eller varmeste Hjærter
i den nyere Tid have gjort gældende for den, som fortræffelige
\emph{Regulativer}, men ikke som mer«. (p. 20). Jeg gad vide, hvad
Etik skulde være andet end netop en Lære om Regulativer, ---
en Lære om deres Udspring og Begrundelse og om de Konsekvenser,
der ved omhyggelig Tænkning og ved Benyttelse af saa mange
Fakta som muligt kunne udledes af dem. Dr. Brandes' Studium
af Etiken synes i flere Henseender ikke at have været heldigt,
siden han ikke en Gang har opdaget dette. Om Etiken end blev
langt mere videnskabelig end den har Udsigt til at blive, saa vilde
den dog aldrig blive andet end en saadan Lære om Regulativer.

Etiken er, tror jeg, i vore Dage kommen ind paa en sundere
og frugtbarere Metode end den i den forudgaaende Periode anvendte;
men det følger af sig selv, at paa dette Omraade ville stedse af%
% --- p. 132 ---
\setcounter{footnote}{0}%
\opage{132}gørende Modsætninger gøre sig gældende. Dr. Brandes gør dog
disse Modsætninger større end de ere, idet han paa aldeles udvortes
Vis river nogle Sætninger fra forskellige Forfattere ud af deres
Sammenhæng og benytter den forskellige Ordlyd til at frembringe
Skinnet af fundamental Uoverensstemmelse. Saaledes stiller han
\emph{Steinthals} Sætning: »Kulturen i og for sig er altid sædelig«
op imod min Sætning: »Ikke al Kulturvirksomhed er etisk«.
Men hvis Læseren vil efterse det af Dr. Brandes citerede Sted hos
Steinthal (Allg. Etik p. 169) i dets hele Sammenhæng, vil han
finde, at Steinthal ved »Kultur« tænker paa al saadan aandelig
og materiel Virksomhed, ved hvilken »den sædelige Virksomhed
fremmes«. Kulturen er efter Steinthal kun sædelig, \emph{fordi} den
fremmer sædelig Virksomhed. Jeg tager Ordet Kultur i videre
Betydning end Steinthal, nemlig om al Bearbejdelse af det naturligt
givne, om enhver Brug af Kræfterne, der fører ud over den Tillæmpelse,
som allerede er vunden. Og derfor siger jeg: Kulturen
 er kun etisk, \emph{for saa vidt} den udvikler de menneskelige Kræfter
 paa en sund og harmonisk Maade. Det er Kulturen (efter min
 Betydning af Ordet), der har hidført det sociale Spørgsmaal, hvis
Forgreninger egentlig alle specielle etiske Spørgsmaal ere. Mit
»\emph{for saa vidt}« rimer sig fortræffelig med Steinthals »fordi«, naar
man bemærker hele Sammenhængen. --- Dr. Brandes siger fremdeles
om Steinthal: »Som radikalere Tænker end Høffding anser han
Ægteskabet, som vi kende det, for uetisk«. Hertil maa jeg bemærke,
at hvad jeg har skildret i min Etik under Navn af »det frie
Monogami« udtrykkelig er betegnet som \emph{et Ideal}, skøndt ikke et
Ideal, der svæver i Luften, da ethvert lykkeligt og alvorligt Ægteskab
indeholder en \emph{Tilnærmelse} dertil (Etik p. 182). Jeg har ikke
lagt Skjul paa den bestaaende Ordnings Ulæmper, og jeg har
fremhævet, at først ved fyldigere Udvikling af Kvindens aandelige
og materielle Selvstændighed kan Tilnærmelsen til det angivne
Ideal blive større (Etik p. 188 f. 195). I Virkeligheden er jeg paa
dette Punkt ikke uenig med Steinthal. Jeg er heller ikke uenig
med ham i hans Opfattelse af Socialismen. Paa ham tænkte jeg
netop, da jeg skrev følgende Linjer: »Dersom Socialismen --- i
Betydning af Fælleseje og Fællesbrug af Arbejds- og Forbrugsmidler
--- udvikler sig som Frugt af fri Sammenslutning, og dersom
den kan bestaa uden en stærk Centralmagt, saa behøves der
intet nærmere Bevis for, at den forudsætter en anden Menneskenatur
end den, vi kende, og den maa altsaa haabe, at der kan
% --- p. 133 ---
\setcounter{footnote}{0}%
\opage{133}foregaa en stor og gennemgribende Forandring af Menneskenaturen.
Dette Haab kunne vi dele« (Etik p. 251). Jeg lægger kun større
Vægt end Steinthal paa Vanskelighederne ved dette Haabs Virkeliggørelse
og paa Nødvendigheden af at finde Udgangspunkter for
Arbejde og Udvikling i den forønskede Retning. Vistnok vilde jeg,
naar jeg nu skulde udarbejde min Etik, anlægge Fremstillingen
saaledes, at jeg kom Socialismen en Streg nærmere end da jeg
for snart syv Aar siden gjorde Udkast til Afsnittet om det sociale.
Spørgsmaal.

Hvad angaar Forholdet mellem \emph{Friederich Paulsens} og
min Etik, kan jeg anføre et Faktum, som ogsaa viser, at Uoverensstemmelserne
mellem forskellige Etikere ikke have saa meget at
sige, som Dr. Brandes paastaar. Jeg haaber ikke at begaa nogen
Indiskretion ved at anføre en Udtalelse af Friederich Paulsen i et
Brev til mig af 2. Juli 1887, efterat jeg havde sendt ham den
danske Udgave af min Etik. Han skriver: »Deres Etik traf mig
i Færd med at udarbejde min egen Etik til Offentliggørelse. Jeg
skyndte mig at læse de Partier, med hvilke jeg netop var beskæftiget,
dem, som behandle de sociale Problemer. Mit Haab om at møde
Tanker, der bares af samme Grundanskuelser og i det væsentlige
førte til samme Løsninger, blev ikke skuffet. Jeg har med stor
Glæde og Tilslutning læst Deres Udviklinger angaaende disse
Punkter. Semestrets Begyndelse afbrød mit Arbejde og min Lekture.
Nu har jeg endelig læst de første Afsnit,* som behandle Principerne.
Overensstemmelsen mellem vor Maade at se og fremstille Tingene
paa var mig ofte overraskende«. --- Jeg kan hertil føje, at i de
Samtaler, jeg tidligere havde haft med med Prof. Paulsen, vare
de herhenhørende Spørgsmaal ikke komne frem. Vi vare altsaa
uafhængigt af hinanden komne til aldeles lignende Opfattelser og
Resultater. Jeg har saa meget des mere Grund til at mindes
Prof. Paulsen, naar jeg slaar et Slag for min Etik, som det var
ved hans velvillige Medvirken, at den tyske Udgave af denne Bog
udkom (nogle Maaneder før hans egen).

Mange af de Uoverensstemmelser mellem moderne Etikere,
Dr. Brandes mener at have fundet, ville saaledes ved opmærksom
Opfattelse falde bort. Men, som sagt, forskellige og modsatte
Retninger ville paa dette Omraade stedse gøre sig gældende. Den
videnskabelige Karakter af den etiske Diskussion vil bero paa, i
hvilken Grad de, der tage Del i den, støtte sig til bestemte Grunde
og til Fakta. Derved udelukkes ikke »Stemningsrigdom«. Det
% --- p. 134 ---
\setcounter{footnote}{0}%
\opage{134}vilde endog være unaturligt at drøfte etiske Spørgsmaal, uden at
det i Aand og Tone mærkedes, at det er store og afgørende Ting,
det drejer sig om. Naar blot Stemningen følger Tanken, i Stedet
for at lede den, er den ingen Hindring for den rationelle Begrundelse.
Men jeg maa gentage min Paastand fra min første Artikel, at
baade Nietzsche og Dr. G. Brandes forveksle en kraftig Stemningsudtalelse
med begrundende Tænkning. Brandes citerer i sin anden
Artikel et Brev fra Nietzsche, i hvilket denne erklærer, at han har
»en Mistro til Dialektik, endog til Grunde«: Udtalelsen (hvis
Interesse fra et rent psykologisk Synspunkt jeg ikke miskender)
synes dog at være lidt underlig, naar den skal tjæne som Bevis
for Mandens filosofiske Interesse. I Filosofien vil der stedse vedblive
at være mange ubekræftede Paastande, mange Hypoteser.
Men en Hypotese er ikke et Paafund, som et begavet Individ har Mod
til at udtale. En Hypotese er en \emph{begrundet} Formodning, en
ved stræng Tænkning vunden Konsekvens, som kun endnu ikke
har fundet, og maaske ikke kan finde sin Bekræftelse ved Erfaring.
Der dukke saa mange Stemninger op, som kunne fortjæne at faa
et klart og kraftigt Udtryk; men kun gennem \emph{alsidig} Prøvelse
kunne de give Stødet til at vinde Klarhed over Virkelighedens og
Livets Forhold. Under en saadan Prøvelse faar man Brug for
mange Sætninger, der fra et. æstetisk Synspunkt kunne synes
»trivielle« eller »næsten-trivielle«, især naar man river dem ud
af deres Forbindelse. »Og« og »men« ere ogsaa trivielle, forslidte
Ord, og dog nødvendige for at udtrykke en Sammenhæng. Det
er lettere at naa en pointeret Form, naar man uden videre skyder
alle Sagens forskellige Sider tilbage med Undtagelse af én eneste.
Men den naaes da ogsaa paa Sandhedens Bekostning. »Fornemmere«
(for at bruge et af den Nietzscheske Skoles Yndlingsudtryk) er det
maaske. Men allerede Kant klagede i sin Tid over en vis »fornem
Tone i Filosofien«, der viste sig i, at man ansaa det for at være
under sin Værdighed at give den alsidige og besindige Begrundelse
af sine Paastande. At denne Klage passer paa Nietzsche, er jeg
ikke ene om at hævde. I en Anmældelse i »Deutsche Rundschau«
(1887), som først nu kommer mig for Øje, skriver Georg v. Gizycki
om et af Nietzsches Skrifter: »Bidrag til virkelig Videnskab har
jeg ikke fundet i Bogen. Dette vil Forfatteren anse for en Kompliment;
ti »»Was ist der wissenschaftliche Mensch?«« Nietzsche
svarer: »»Zunächst eine unvornehme Art von Mensch««; og han
elsker fremfor alt det fornemme«. ---

% --- p. 135 ---
\setcounter{footnote}{0}%
\opage{135}Jeg tror derfor ikke, at jeg har gjort Nietzsche Uret. Derimod
gør Dr. Brandes mig Uret, naar han anker over, at jeg i
Slutningsbemærkningen af min Artikel »Demokratisk Radikalisme«
har udtalt en Beklagelse over den Sygdom, der har rammet
Nietzsche. Han mener, at denne Beklagelse ikke »klinger rigtig
ægte«. Det kommer an paa Gehøret. Jeg har ikke sagt andet,
end hvad jeg virkelig føler overfor en Mand, jeg selv har anerkendt
som en fremragende Forfatter. Og min Bemærkning refererede
sig til et Sted i Dr. Brandes' \emph{egen} Artikel, hvor han omtaler
Nietzsches Sindssygdom! Kun én Ting tog jeg ikke med fra dette
Sted hos Dr. Brandes: nemlig hans Forklaring af Sindssygdommens
Udspring af potenseret Selvfølelse. Om der i den offentlige
Forkyndelse af denne Forklaring ligger nogen særlig Hensynsfuldhed
fra Dr. Brandes' Side overfor den af ham højt beundrede
Forfatter, til hvem han har staaet i et vist personligt Forhold,
vil jeg lade staa hen. Hvor det gælder Takt, ere Meningerne jo
meget forskellige.

--- En ikke ringe Del af sit Svar anvender Dr. Brandes til
et Forsøg paa at godtgøre, at jeg i forskellige Henseender har
misforstaaet Nietzsche. Jeg skal gøre Rede for, hvorledes det dermed
forholder sig.

1. Inden jeg i »Demokratisk Radikalisme« gik over til at
belyse de enkelte Paastande i Dr. Brandes' første Artikel, havde
jeg i Forbigaaende henvist til den ovenfor omtalte karakteristiske
Overensstemmelse mellem Nietzsche og Brandes, at de ere tilbøjelige
til at overse Forskellen mellem en Overbevisnings Styrke
og dens Gyldighed. Jeg anførte et Par Steder af Nietzsches i
filosofisk Henseende vigtigste Værk, i hvilket han taler imod den
videnskabelige Nøjagtighed og Ængstelighed som Rest af Askesen
og Dogmatismen. Dr. Brandes giver mig en hel Lektion i hele
dette Afsnit hos Nietzsche, en Lektion, jeg aldeles ikke havde
behøvet, da jeg har forstaaet det ganske paa den Maade, han
udvikler. (Kun én Indrømmelse maa jeg gøre. Jeg har ikke gengivet
Nietzsches Udtalelse ordret. Jeg har ladet ham sige: »Hvis
nu intet var mere guddommeligt end Løgnen!« i Stedet for:
»Hvis nu intet mere var guddommeligt, det skulle da være Løgnen!«
Denne Forskel gør ikke stort til Sagen, skøndt Dr. Brandes gør et
Nummer af den. Nietzsche udtrykker sig blot noget \emph{stærkere}
end i min Gengivelse). Den hele Tankegang hos Nietzsche er
aldeles ikke ny. Den gaar i Korthed (for at benytte Dr. Brandes'
% --- p. 136 ---
\setcounter{footnote}{0}%
\opage{136}Resumé) ud paa, at »da i Livet baade Sandhed og Usandhed
stadigt vise sig nyttige, kan den videnskabelige Overbevisning [om
Sandhedens Vigtighed] ikke have sin Oprindelse i en Beregning af
Nytten, men maa være opstaaet til Trods for, at Unytten og
Faren ved at ville Sandhed \emph{for enhver Pris}, stadigt læres af
Livet.« Hvorfor ikke bedrage og lade sig bedrage? --- Den Tankegang,
som hermed er antydet, er under forskellige Former kommen
frem i Sammenhæng med den kritiske Filosofi, Udviklingslæren
og Velfærdsmoralen, sjældent dog i en saa forskruet Form som
hos Nietzsche. Hans Fortjæneste er her kun at overdrive den og
dernæst indklæde den i sin sædvanlige Profetstil. --- Problemet
er: hvorledes forholder Erkendelsesdriften sig til Kampen for Tilværelsen?
Da det ofte er et Spørgsmaal paa Liv og Død, om
der sanses og tænkes ret, vil det være let at forklare, hvorledes
Erkendelsesdriften kunde vækkes under Kampen for at bestaa.
De Racer, der aldeles ikke havde nogen Trang til Sansning og
Tænkning, eller ikke tilstrækkelig Trang, vilde gaa under, da de
vilde blive overfløjede og fortrængte af Racer med skarpe Sanser
og energisk Tænkeævne. Ganske vist, dersom Nietzsche havde
Ret i, at Usandhed spillede lige saa stor en Rolle i Tilværelseskampen
som Sandhed, vilde Erkendelsesdriften være uforklarlig.
Kun kunde man sige, at Ævnen til at lyve og bedrage forudsætter
Ævne til at finde Sandheden; ti Bedrageren maa kende de virkelige
Forhold for at kunne anvende sin List paa rette Maade. Det er
nu en tydelig Overdrivelse, Nietzsche gør sig skyldig i, naar han
antyder, at en Drift til Usandhed og Bedrag lige saa godt kunde
have udviklet sig, som den modsatte Drift. Kampen for Tilværelsen
nøder det menneskelige Samfund til at trænge Løgnen
tilbage blandt dets Medlemmer indbyrdes. Gensidig Aabenhed og
Sanddruhed er en Livsbetingelse, som kun til en vis Grad kan
krænkes uden fuldstændig Opløsning af alle Samfundsforhold.
Naar Driften til Erkendelse og Sandhed nu til en vis Grad begunstiges,
og begunstiges mere end den modsatte Drift, saa vil
ogsaa her Motivforskydningen gøre sig gældende: hvad der først
øvedes som Middel i Kampen, vil øves for dets egen Skyld, med
umiddelbart Velbehag ved selve Kræfternes Brug, afset fra det
praktiske Udbytte. Der danner sig en ren intellektuel Drift, og
den intellektuelle Glæde bliver mulig. Dette sker i Kraft af den
almindelige psykologiske og fysiologiske Lov, at naar en Virksomhed
først er bleven tilstrækkelig indøvet, vil den kunne foregaa
% --- p. 137 ---
\setcounter{footnote}{0}%
\opage{137}med Lethed og paa uvilkaarlig Vis, selv om den først krævede
Anstrængelse og maatte udløses ved ydre Paavirkning. Saaledes
opstaar den videnskabelige og den filosofiske Interesse, som forfølger
sit Maal, selv om intet praktisk Udbytte er i Sigte, og selv
om det viser sig haabløst at udtømme og omspænde Virkelighedens
uendelige Indhold. En videre gaaende Betragtning vil dog
vise, at alle videnskabelige Sandheder og alle filosofiske Ideer af
Betydning ikke skyldes den blotte intellektuelle Interesse, men at
Tænkningen, det være ad nok saa mange Omveje, tilsidst stedse
godtgør sin Betydning for den praktiske Udvikling af Menneskelivet.
Hvor Tænkningen er kommen for langt bort fra Livet, er
den gærne i Færd med ogsaa at komme bort fra, hvad der tjæner
dens egen Udvikling.

Jeg tror at kunne sige, at hvis Nietzsche havde set mere
ædrueligt og alvorligt paa Tingene, vilde han ikke have spildt
saa megen uægte Dybsindighed paa dette Punkt. Havde han haft
mere virkelig filosofisk Interesse, vilde hans Stemningsleg paa
Filosofiens Grænse ogsaa haft større Værdi.

2. Mod min Bemærkning om, at Nietzsche ikke har frigjort
sig fra Schopenhauer, strider det ikke, at han forskellige
Steder polemiserer mod denne. Som jeg har paavist, benytter han
stærkt den Schopenhauerske Tohed: Væmmelse ved Livet ---
Medlidenhed med Livet. Men medens Schopenhauer stiller Medlidenheden
i Modsætning til Væmmelsen og overvinder denne ved
hin, bliver den Medlidenhed, der findes i Verden, for Nietzsche
kun en saa meget des større Grund til at væmmes ved Verden.
I denne Forstand kan man sige, at han stadigt bruger de Schopenhauerske
Briller. Havde han været helt frigjort for Schopenhauer,
havde hans Tanker modtaget kraftige Impulser fra helt andre
Sider, saa vilde hele hans Opfattelse være kommen 'til at se
anderledes ud. Gennem Schopenhauer hænger han tilbage sammen
med Romantiken i Aarhundredets første Decennier.

3. Overfor min Paastand om, at Nietzsche forveksler Spørgsmaalet
om Moralens \emph{historiske Oprindelse} med Spørgsmaalet
om dens \emph{Værdi}, anfører Dr. Brandes flere Steder, hvor Nietzsche
omtaler Forskellen mellem disse to Spørgsmaal. Saadanne Steder
kendte jeg meget vel. Men det for mig afgørende var, at Nietzsche
i det Skrift, i hvilket han har fremstillet sin i filosofisk Henseende
betydningsfuldeste Teori, næmlig »Genealogie der Moral«, afgør
Spørgsmaalet om Moralens Oprindelse \emph{som om} Spørgsmaalet om
% --- p. 138 ---
\setcounter{footnote}{0}%
\opage{138}dens Værdi dermed ogsaa var afgjort. Han erklærer udtrykkeligt
Behandlingen af det første Spørgsmaal for Middel til Afgørelsen
af det andet (Genealogie der Moral, p. IX), og begge Dele slaa
hele Bogen igennem over i hinanden. Af sin Teori om Moralens
Oprindelse udleder han en Forkastelse af den hidtil udviklede
Morals Værdi, og paa den anden Side er det tydeligt nok, at han
har lagt sin egen sygelige Antipati mod Demokrati og Socialisme
over i de gamle Slægters Tankegang, i hvis Foragt for de lavere
staaende Stænder eller Racer han .finder Moralens Udspring.

 4. Hvad selve Nietzsches »Moralgenealogi« angaar, fandt
jeg en Modsigelse i den, da han anser de reaktive Følelser, d. v. s.
dem, der fremkaldes ved Trang til Modvirken overfor andre, for
slettere end de aktive, som ikke i den Grad ere bestemte udefra.
»Det aktive, angribende, overgribende Menneske«, siger Nietzsche
(Geneal. der Moral p. 64), »staar bestandig Retfærdigheden hundrede
Skridt nærmere end det reaktive; det er for ham aldeles ikke
nødvendigt at vurdere sin Genstand paa falsk og forud indtagen
Vis, saaledes som det reaktive Menneske gør og maa gøre«. Jeg
bemærkede hertil, at Foragt lige saa vel er en \emph{reaktiv} Følelse
som Nag, og at jeg ikke kunde indse, hvorfor den forulæmpedes
og nedtrampedes Hævnlyst skulde være lavere og uædlere end
Selvfølelsen og Foragten hos den, der tramper paa ham. Foragt
og Nag svare i psykologisk Henseende ganske til hinanden. Begge
opstaa kun ved bestemt udprægede Modsætningsforhold og regnes
derfor gærne af Psykologerne til hvad man paa Tysk kalder
»Erwiederungsgefühle« (Smlgn. A. \emph{Horwicz}: Psychologische Analysen.
II, 2. p. 325 f.). I sin aristokratiske Iver har Nietzsche her
gjort Vold paa Psykologien.

Dr. Brandes mener, at den aristokratiske Selvfølelse og
Foragt overfor andre vel er en \emph{relativ} Følelse, men ikke nogen
\emph{reaktiv} Følelse. Denne Distinktion kender Nietzsche ikke. Lige
efter at han, som anført, har erklæret det »aktive« Menneske for
bedre end det »reaktive«, bebrejder han den moderne Biologi, at
den har borteskamoteret Begrebet om den egentlige Aktivitet.
»Man stiller »»Tillæmpelsen«« i Forgrunden, det vil sige, en Aktivitet
af anden Rang, en blot Reaktivitet« (Anf. Skr. p. 69). Altsaa
han hævder udtrykkeligt, at der skal gives Følelser, som ikke bero
paa Tillæmpelse, d. v. s., ikke ere bestemte ved Forhold til noget
ydre, ikke ere relative, og at den aristokratiske Selvfølelse skal
være af denne Art. Har jeg da ikke Ret i min Indvending, at
% --- p. 139 ---
\setcounter{footnote}{0}%
\opage{139}en saadan Selvfølelse ikke \emph{kan} opstaa uden i Modsætning til og
som aandelig Modvirken mod Svagheden og Usselheden, og at den
altsaa lige saa vel er en »reaktiv« Følelse som Naget, der udspringer
som Modvirken mod Forulæmpelse? --- Dersom »reaktive« Følelser
ere af lav Natur, saa passer det paa Foragten lige saa vel som
paa Naget, og hans Teori om hins naturlige Ædelhed falder
sammen. ---

Ved sin Lovprisning af Foragten og ved sin Paastand om
Naget og Misundelsen som herskende hos den store Masse viser
Nietzsche sig besjælet af en af de værste af de Stemninger, som
Militarismen har fremkaldt indenfor det tyske Folk i visse Kredse.
Overmod og Ringeagt have udviklet sig hos de herskende Klasser,
naturligvis med mange hæderlige Undtagelser. Selv ved de Statsforanstaltninger,
man hidtil har sat i Værk for at lindre de sociale
Ulæmper, ytrer der sig baade i Anlægget og Gennemførelsen en
Mistanke mod enhver \emph{selvstændig} Livsytring fra neden, enhver
Fordring om Selvstyrelse. Man er tilbøjelig til at se den hele
socialdemokratiske Bevægelse som udsprungen af Nag og Misundelse,
hvad den dog sikkert \emph{ikke} er i nogen højere Grad end en hvilken
som helst anden Standsrejsning, om hvilken Historien bærer
Vidnesbyrd. Noget af det, som er Nietzsche mest imod i Nutiden,
er dens Utilbøjelighed til at lade sig regere, hvad han kalder dens
»Misarchisme«. Det synes da, at den Selvraadighed, Dr. Brandes
i Tilslutning til Nietzsche har anbefalet, i Følge Mesteren selv
allerede praktiseres i fuldt Maal, --- rigtignok af de mange, ikke
af de udvalgte faa. Men hvorledes skal man drage Grænsen?
Og var det i hvert Tilfælde ikke bedst ikke at \emph{tale for meget
om den nye Lære}, for at dens Praktisering kunde blive forbeholdt
de udvalgte? ---

Det vilde efter den literaturhistoriske Metode, som Dr. Brandes
plejer at anvende (rigtignok sjældent paa Forfattere, han i Øjeblikket
sympatiserer med), have ligget ham nær at betragte Nietzsche
som Udtryk for visse Strømninger indenfor det moderne tyske
Aandsliv. Baade paa Grund af Nietzsches Begavelse og paa Grund
af hans Sygelighed faar saa meget i den almindelige Tilstand Luft
gennem ham paa særlig intensiv Maade. I Lyset af sine Stemninger
betragter han saa, hele Verdenshistorien, især Moralens og Religionens
Udvikling. Ikke som filosofisk Teori, men som Replik
eller Monolog have hans Ord Interesse. ---
% --- p. 140 ---
\setcounter{footnote}{0}%
\opage{140} Inden jeg forlader Nietzsche, vil jeg endnu slaa fast, at
Dr. Brandes i sin anden Artikel aldeles \emph{ikke har svaret noget
som helst paa selve Hovedpunkterne i min Kritik af
Nietzsches Moralgenealogi} (se min Artikel »Demokratisk
Radikalisme«, »Tilskueren«. November-December Hæftet 1889,
p. 869 f. og 871 f.), altsaa \emph{af det hos Nietzsche, hvorpaa
den hele Teori om aristokratisk Radikalisme tilsidst
grunder sig}. --- Dette vil vise sig at hænge sammen med den
ejendommelige Beskaffenhed af Dr. Brandes' Svar i det hele.

\sectionnum{III.}

Det var to Hovedpunkter i Artiklen »Aristokratisk Radikalisme«,
min Indsigelse drejede sig om: 1) Teorien om de store Mennesker,
og 2) Velfærdsmoralens Gyldighed. Lad os se, hvorledes disse to
Stridspunkter stille sig efter Dr. G. Brandes' sidste Udtalelse.

1. Hvad der gav Dr. Brandes' Artikel »Aristokratisk Radikalisme«
dens Betydning og voldte dens store Virkning, var, at
naar man holdt sig til den hele Stemning og Tankegang i den
og særligt mærkede sig enkelte Kraftsteder, maatte man faa ud,
at han (i Tilslutning til Nietzsche) proklamerede, at den historiske
Udviklings Maal var naaet i de forholdsvis faa store Mennesker, i
Forhold til hvilke alle andre kun vare Midler og Omveje, naar de
da ikke vare Hindringer. Det er denne Sætning, som --- hvis
den virkelig ligger i Artiklen --- giver denne dens Nyhed og berettiger
den Opsigt, den har vakt. Det er denne Sætning, og kun
den, jeg har angrebet. Men jeg tog straks et Forbehold, idet jeg
ikke klart kunde se, om det virkelig var Meningen at opstille
den nævnte Sætning, da andre Udtalelser tydede paa, at hine
store Mennesker dog \emph{tillige} skulde betragtes som Midler og Redskaber
for Slægtens \emph{fortsatte} Liv og Udvikling. Hvis dette er
Meningen, sagde jeg, bortfalder al min Betænkelighed og hele
min Kritik.

Dr. Brandes har nu selv erklæret, at dette sidste har været
hans Mening, og Sagen er for saa vidt klar. Han maa jo selv
bedst vide, hvad han har ment. Jeg vil kun sige, at han da har
været meget uheldig med at udtrykke sig. Hans Erklæring vil
sikkert vække Forbavselse hos Læserne af hans første Artikel,
ikke mindst hos dem, der i god Tro have dateret en ny Æra fra
de store Ord (ti andet var det altsaa ikke), den indeholdt.

% --- p. 141 ---
\setcounter{footnote}{0}%
\opage{141}\emph{Faktisk har Dr. Brandes ved sin Erklæring brækket
Spidsen af de store Ord og de dristige Tanker i sin
første Artikel}. Det Hovedindtryk, jeg for mit Vedkommende
har faaet af hans anden Artikel, har mindet mig om den brave
Præst, der, da han saa Menigheden saa stærkt bevæget ved hans
Præken, trøstede den med disse berømte Ord: »Græd ikke, lille
Børn, det turde være Løgn alt sammen!«

. --- Dog giver denne anden Artikel mig, hvad Spørgsmaalet
om Stormennesker angaar, Anledning til nogle Bemærkninger.

Der bebrejdes mig, at jeg gør en »plump, udialektisk Forskel
mellem Maal og Middel«. Jeg føler mig ikke ramt herved.
Jeg har benægtet, at de store Mennesker \emph{kun} vare \emph{Maal, de}
andre kun Midler. Den Opfattelse, \emph{jeg} i min Etik havde søgt at
gennemføre, var netop den, at ethvert Væsen med \emph{Ævne til at}
føle Lyst og Smerte er at betragte \emph{paa én} Gang som Maal og
Middel. Dette Princip afgiver en Maalestok ved Undersøgelsen af
de enkelte Livsforhold. Naturligvis vil Forholdet mellem Maal og
Middel være højst forskelligt for de forskellige Væsner efter deres
indre og ydre Vilkaar. Jeg ser Sagen paa den Maade, at ved
Ordet Maal betegnes her særligt den passive \emph{Side, den} Side, fra
hvilken vi forholde os modtagende. Barnet er saaledes Maal
for Moderens Omhu og Arbejde. Ved Ordet Middel betegnes derimod
her den aktive Side, den Side, fra hvilken vi forholde os
handlende. Derfor fandt jeg i min første Artikel det mest
karakteristiske for de store Mennesker i, at de ere \emph{de store Midler,
de store Tjænere}. De have jo Ævne til ved deres Væren og
Virken at blive \emph{Udgangspunkter} for stærke og varige Strømninger,
der bringe aandelige og fysiske Livsbetingelser og Livsspirer
til Veje for store Kredse. Som Goethe har udtalt det: de
store Geniers Produktivitet strækker sig ud over deres egen Person
og Tid og virker maaske Aarhundreder igennem. Derfor kunde
jeg ikke gaa ind paa den »simple Tro til, at Menneskehedens
Maal ikke er et i det \emph{uendelige udskudt}, men maa ligge i
dens højeste Eksemplarer«. (Brandes: Aristokratisk Radikalisme,
p. 575). Jeg vilde netop sige: den største er den, der sætter
noget i Bevægelse, hvis Virkninger gaa i det uendelige, og som
aldrig ophører at virke som Middel til Liv og Lykke for de mange.
Ordet \emph{sidste Maal} synes dog uheldigt valgt for det, hvis Betydning
netop er at være Udgangspunkt. Kraften kende vi kun gennem
dens Virkninger, og efter dens Virkninger 'maa vi vurdere den.

% --- p. 142 ---
\setcounter{footnote}{0}%
\opage{142}Jeg har ingenlunde lagt Vægten paa, hvor meget det store
Menneske selv \emph{føler sig} som Middel. Jeg har kun i Forbigaaende
sagt, at de store Mennesker selv ganske sikkert vilde afvise den
Nietzscheske Teori, fordi de hverken i deres Modtagen eller i deres
Virken føle sig sondrede fra den øvrige Menneskehed, og fordi de
ene i den Idé, i den Sag, de tjæne, ville finde Maalestokken for
deres Værd. At Linen i dem skulde være løben ud, ville de ikke
gaa ind paa. Jeg har udtrykkelig gjort opmærksom paa, at
Stormennesket hverken behøver \emph{at tænke paa} Menneskehedens
Velfærd eller en Gang at have nogen klar Bevidsthed om sin
Opgaves Betydning. Det største øves ofte uvilkaarligt, drømmende,
instinktmæssigt. Det Motiv, af hvilket vi erklære et Menneske for
stort, \emph{behøver} ikke at være det, der har besjælet dette Menneske
selv ved dets Handlen. Eller som jeg plejer at udtrykke det:
Vurderingsmotiv og Handlingsmotiv behøve ikke at falde sammen.
Og det kan ofte være heldigt, om de ikke gøre det. Naar en stor
Mand falder i Vandet, vil det næppe være heldigt, om han vil
give sig til at beregne, om det nu ogsaa for Menneskehedens Vel
er nødvendigt, at han kommer op. Derimod vil det være heldigt,
om hans Selvopholdelsesinstinkt er saa stærkt, at han uden videre
arbejder sig op. Og saaledes forholder det sig ogsaa med al
Virken i det store, ja i Grunden med al Virken. Det sidste Maal
tjæner ved de enkelte Handlinger og Værker kun som en Maalestok,
et Kriterium, der i Regelen kun højst indirekte og gennem
mange Mellemled, som langvarig Refleksion maa udfinde, kan
anvendes paa de særegne Tilfælde. Baade Tankeløshed og Hensynsløshed
kunne derfor ofte være at anbefale.

Ogsaa jeg betragter det store Menneske som Kulturens
Kilde. Det udtalte jeg i »Demokratisk Radikalisme« p. 855 f. 857 f.,
og det var allerede en Hovedtanke i min Etik, hvor den navnlig
bliver af Betydning ved Spørgsmaalet om Statens Forhold til Kulturen.
Jeg har dér hævdet, at Kulturen ikke er Statens Værk, ti
»baade den materielle og den aandelige Kulturs Oprindelse fører
os stedse tilbage til enkelte Individer, i hvis Bevidsthed de nye
Tanker ere blevne til. Indenfor den enkeltes Bevidsthed have
disse nye Tanker deres første Kamp for Tilværelsen at bestaa\dots{}.
Hvad der saaledes har udviklet sig i det skjulte, kan senere
komme Statslivet til Gode, men dette kan ikke selv fra først af
udvikle det« (Etik p. 366).

% --- p. 143 ---
\setcounter{footnote}{0}%
\opage{143}\emph{Store Mennesker kunne altsaa aldrig undværes}.
Dr. Brandes mener nu, at jeg siger mig selv imod, naar jeg
(overfor Overvurderingen og Forgudelsen af de store Mennesker)
udtaler, at \emph{det dog ofte var bedst, om de havde kunnet}
undværes. Hvilken Selvmodsigelse ligger der heri? Det er jo
dog en klar og simpel Tanke, at et stort Menneskes historiske
Uundværlighed ofte kan bero paa, at Udviklingen har taget en
ulykkelig Vending, at man er kørt fast, at Samfundet er ved at
opløse sig. Livet er blevet tomt, betynget, delt. Da kommer den
store som Frelseren. Men var det dog ikke bedre, om de mange
havde kunnet rejse sig ved egen Hjælp? Det er jo dog det,
Hjælperen vil: give dem Kraft til at hjælpe sig selv! Mig forekommer
Selvmodsigelsen netop at være paa deres Side, der mene,
at den store Hjælper er Formaalet, og at han fremfor alt maa
komme til at anvende sin Ævne til at hjælpe, hvilket forudsætter,
at man maa bringe Folk i Ulykke, for at denne Ævne ikke skal
ligge ubrugt hen! Det skulde altsaa ikke have været heldigst for
Frankrig, om det ikke havde behøvet Jeanne d'Arc til at hjælpe
sig? Eller naar en genial Advokat ved sin glimrende Skarpsindighed
og Veltalenhed frelser et Menneske, der var uskyldigt anklaget af
lumske Fjender, saa skulde det ikke være bedre, om der ingen
lumske Fjender og kortsynede eller partiske Dommere var? ---
Det forekommer mig at være Livets egen Ret, det her gælder,
og det forekommer mig, at jeg med Grund har brugt Ordet
Moloksdyrkelse om den Anskuelse, jeg her bekæmper.

Man kan sige: det er jo dog kun et fromt Ønske, at de
store Mennesker \emph{skulde kunne} undværes, naar de nu faktisk
\emph{ikke kunne} undværes. Jeg har heller ikke fremført det som
andet. Men det tjæner dog til at belyse det hele Spørgsmaal.
Større Vægt lægger jeg ikke paa denne Betragtning. Jeg forstaar
ogsaa meget godt den Følelse, der i den gamle Kirke førte til
Udraabet: felix culpa! d. v. s. »Godt, at vi have syndet, --- ellers
havde vi ikke faaet en saadan Frelser!« Det er den samme
Tankegang, som gør,* at Clärchen priser Kærlighedens Lykke, der bestaar baade
i at \emph{sørge} og juble. (Smlgn. »Demokratisk Radikalisme«
p. 861 f.). I det virkelige Liv ere Sorg og Lyst, Nød og Hjælp,
Opløsning og Udvikling, saa inderligt knyttede til hinanden, at vi
ikke kunne tænke os, hvorledes de menneskelige Kræfter skulde
kunne udvikles og anvendes uden netop gennem en saadan Bølgegang.
Det bliver altid kunstigt, naar vi tænke det ene Modsæt%
% --- p. 144 ---
\setcounter{footnote}{0}%
\opage{144}ningsled bort. Men med mit Tankeeksperiment har jeg netop
villet vise, hvor urimelig den modsatte Tanke er, der lader de
store Mennesker være et og alt (eller tenderer dertil).

De to Opgaver: \emph{at stifte ny Kultur} og \emph{at udbrede
den vundne Kultur}, behøve i Virkeligheden ikke at komme i
Strid med hinanden. Der er da ingen Grund \emph{til} at tage ensidigt
Parti for nogen af dem, og det er unaturligt at stille dem i Modsætning
til hinanden. Dertil kommer, at --- som allerede bemærket,
--- det store, det kulturstiftende oftest arbejder sig frem
i de udvalgte Menneskers Indre paa en overvejende instinktmæssig
Maade. Først naar det er kommet til en vis Udvikling ad denne
\emph{uvilkaarlige Vej, kan} det blive Genstand \emph{for besindig og} vilkaarlig Bearbejdelse, Tilegnelse og Udbredelse. Derfor kunne
ydre Anstrængelser, særlig Statsforanstaltninger gøre saa overordentligt
lidt for at stifte nye Kulturstrømninger. De saakaldte
»Kulturlande« kunne ofte, være længst fra at stifte ny Kultur.
Det mest storartede, som er kommet frem paa Aandslivets Omraade
i den nyeste Tid, den moderne russiske Literatur, kom fra
et Folk, man i Vesteuropa havde vænnet sig til at betragte som
barbarisk. Ingen Statsforanstaltning vilde her kunne hjælpe.

Der findes i Dr. Brandes' anden Artikel en noget uklar
Udtalelse. »I Reglen,« siger han (p. 24), »føles ikke Spaltningen
mellem den demokratiske og den aristokratiske Betragtningsmaade,
fordi man samtidigt i Staterne gør adskilligt for Massernes elementære
Kultur og en lille Smule for de kulturbærende Elementers
Vel.« Menes der med »kulturbærende« Elementer saadanne,
som stifte ny Kultur, da staar Staten afmægtig overfor
dem. Kun højst indirekte kunne Statsforanstaltninger hjælpe her.
\emph{Ingen Statsmand kan frembringe »Kulturkilder.« De
Lag, fra hvilke disse Kilder komme, ere endnu} ikke
\emph{opdagede. Netop fordi man ikke ved, hvor disse Lag
ligge, er den rigtigste Vej at gaa den: at skaffe saa
gunstige Livsbetingelser for saa mange som muligt}.
»Derved forøges,« som jeg har vist i »Demokratisk Radikalisme«
(p. 854—858), »ikke blot Muligheden for, at der kan opstaa
Stormennesker, men ogsaa for, at de kunne finde den Anerkendelse,
Tilslutning og Modvægt, uden hvilke deres Virken ikke bliver
frugtbar og gavnlig.« Skal man sætte \emph{en Opgave for Samfundet}, kan der aabenbart kun sættes: at udbrede den Kultur,
 som er vunden, --- naturligvis ikke blot den »elementære«, men
% --- p. 145 ---
\setcounter{footnote}{0}%
\opage{145}al hidtil vunden Kultur, --- og ikke blot til én, men til alle
Klasser i Folket, eller til saa mange Klasser som muligt. Derved
vil Jordbunden beredes saaledes, at de Spirer, den maatte indeslutte,
ville kunne udfolde sig. I min Etik (p. 367) har jeg set
Statsmandens Genialitet og Betydning deri, at han formaar at
»opdage de nye Livsspirer i Folket og --- direkte eller indirekte ---
yde dem den Hjælp og Støtte, de behøve,« og »at han paa
ejendommelig Maade fortætter aandelige Elementer og Kræfter,
som i \emph{Forvejen} røre sig.«  Men Betingelsen er netop, at de i
Forvejen røre sig.

\emph{Selve Udbredelsen} af \emph{den vundne} Kultur \emph{kan
i øvrigt kræve Stormennesker}, lige saa vel som Opdagelsen
af de nye Kulturspirer. Endnu er det Spørgsmaal ikke løst, i
hvilken Form og paa hvilken Maade Kulturindholdet bedst bibringes
dem, der leve i de ringeste Kaar og have de tarveligste
Forudsætninger. Meget taler for, at der maa finde en hel Omsmæltning,
for saa vidt en Nydannelse Sted, før den folkelige
Tilegnelse af Kulturindholdet kan finde Sted. Rousseau og Pestalozzi
vare Stormennesker lige saa vel som mange, der nævnes
højt i Videnskab, Kunst og Statsliv, og der er endnu Plads for
megen Storhed paa den Vej, de have betraadt.

Der synes ikke at være nogen Grund til at klage over
»Vulgariseringen«. Det virkeligt store vil nok kunne arbejde sig
frem, selv om der anvendes endnu langt mere end hidtil paa
»Massernes elementære Kultur.« Dersom »Vulgariseringen« skulde
føre til, at det store ikke saa let trives og ikke saa let forstaas,
saa er det et Tegn paa, ikke at vi have for megen »Vulgarisering«,
men at den ikke drives paa den rette Maade, --- at ingen
Stormennesker have arbejdet i dens Tjæneste!

Naar det er umuligt at angive en Metode til ligefrem at
skabe Stormennesker, saa er det ikke mindre umuligt at angive
en aldeles sikker Ordning, der kan føre til, at Stormenneskene
altid faa den Indflydelse i det menneskelige Samfund, som tilkommer
dem. Og det af den Grund, at \emph{Striden i den menneskelige
Verden} paa alle Omraader netop staar om, hvad
Storhed er, og hvem der er stor. At sige: Stormenneskene
skulle regere! er et meget barnligt Program, som minder om den
Konstitution, Sismondi udkastede som Dreng, og hvis første Bestemmelser
vare disse: § 1. Tous les Français seront heureux.
§ 2. Tous les Français seront vertueux. At finde en saadan
% --- p. 146 ---
\setcounter{footnote}{0}%
\opage{146}Maalestok for Storheden, at den straks anerkendtes af alle og
sattes paa Herskersædet, vilde være at finde de vises Sten. Den
er hidtil kun funden i Drømmenes og Utopiernes Verden. I den
virkelige Verden er den eneste mulige Fremgangsmaade den, at
de store, i det hele de kyndige og erfarne, maa stræbe at overbevise
de andre. Konger, Ministre, Folkemængder --- alle Slags
Magthavere høre til disse »andre«. Kun paa den Maade kan
Sandhed og Ret blive forbundne med Magt.

\emph{Demokrati}, siger Dr. Brandes, betyder Folkeherredømme
 og intet andet. Hertil er at sige, at intet fornuftigt Menneske nu
til Dags ved Demokrati forstaar hele Folkets \emph{umiddelbare}
Herredømme. Herredømmet maa altid udøves af nogle faa.
Spørgsmaalet er, hvorledes disse faa skulle udfindes; de skulde
 jo helst være saa store som muligt. Der synes kun at være tre
Maader, paa hvilke faa kunne herske over mange. Det kan for
det første ske ved fysisk Tvang og ved den \emph{Frygt}, denne indgyder.
Her forudsættes da, at de faa raade over Vaaben, der
gøre dem det muligt at faa Bugt med de mange. Det kan for
det andet ske ved \emph{uvilkaarlig, pietetsfuld eller beundrende
Tilslutning og Efterligning} fra de manges Side.
Det er Autoritetsforholdet i dets ædleste Skikkelse. De mange
optræde her mere aktivt, end hvor Tvang og Frygt ere det afgørende.
Denne Aktivitet bliver endnu større, naar der fra begge
Sider føles Trang til at konstatere i bestemt og udtrykkelig Form,
\emph{hvor mange} der ville give de ledende deres Tilslutning. Derved
komme vi til den tredje Form, som nu er den almindelige i
civiliserede Lande, og som bestaar i, at \emph{Folkets gennem Afstemning
udtalte Vilje} giver Ledelsen af de offentlige Anliggender
den sidste Sanktion.

Dersom Dr. Brandes har udfundet en Maade, paa hvilken
et Mindretal kan komme til at herske, og som er forskellig fra
disse tre, burde han ikke for længe lade Menneskeheden i Uvidenhed
om denne Opdagelse. Det kan dog ikke være hans Mening,
at den skal være de lykkelige faas Særeje.

Det skyldes en fuldstændig Forvrængning af mine Udtalelser
i min »Etik«, naar Dr. Brandes (p. 6 f.) lader mig lovprise Flertalsherredømmet i og for sig. Paa den Side i min Etik
(p. 397), som han citerer, minder jeg udtrykkeligt paany om, at
»alt stort og godt er udgaaet fra Enkeltmænd og har fundet sin
første Pleje i mindre Kredse.« Men naar en Ide skal træde ud i
% --- p. 147 ---
\setcounter{footnote}{0}%
\opage{147}Virkeligheden, behøver den Tilslutning, maa vinde Stemningen
for sig. Jo mere aandelig Udvikling og Selvstændighedsfølelse
udbreder sig, des mindre bliver det muligt at gennemføre betydelige
Foranstaltninger ved Tvang eller ved blot at støtte sig til
den dunkle Følelse, en Autoritet kan sætte i Bevægelse. Der
maa gøres en bestemt Prøve paa, om Flertallet er overbevist.
Denne Prøve kan man en Tid lang unddrage sig, naar man
raader over den fysiske Magt; men man vil da ogsaa blive afmægtig
netop paa de væsentligste Omraader, --- og i Længden
vil Flertallets Sanktion vise sig uundværlig. Jeg har ikke sagt,
at Flertallet altid har Ret. Men saa længe Flertallet ikke er
overbevist, saa længe er Diskussionen om offentlige Afgørelser
ikke sluttet. Hvis der nu \emph{maa træffes en Afgørelse}, skøndt
Spørgsmaalet endnu staar aabent, \emph{saa det ikke vides, om
Flertal eller Mindretal har Ret}, saa er det jeg siger, at
det er \emph{rimeligere}, at Flertallets Vilje staar ved Magt Det er
rimeligere allerede af den Grund, at Flertallet vil kunne \emph{sætte}
sin Vilje igennem med Magt (naar da Mindretallet ikke raader
over særegne fysiske Magtmidler, der opveje Tallet). Ved ethvert
Forsøg paa at indrette en Forfatning, vil det, om end ad nok
saa mange Omveje, tilsidst være Flertallets Stemning, det
kommer an paa. Jeg har ikke lagt Skjul paa, at der heraf kan
udspringe tragiske Konflikter, idet Minoriteten lige saa vel kan
have Retten paa sin Side. Under \emph{enhver} Statsforfatning er der
Mulighed for en Konflikt mellem den ideelle og den reelle Magt
Derfor har jeg sagt, at enhver Stat er en Nødstat. Jeg er ingen
ubetinget Beundrer af Parlamentarismen, men har selv fremhævet,
at den i Europa er en historisk betinget Statsform, der har sine
store Ulæmper. (Smlgn. min Afhandling om Alexander Hamilton
i »Tilskueren«s Septemberhæfte 1889). I Afsnittet om Staten i
min Etik har jeg søgt at afveje Forholdenes og Retningernes forskellige
Sider saa omhyggeligt, som det var mig muligt indenfor
den Ramme, jeg havde sat mig. En af mine nærmeste Venner,
der ikke er enig med mig i politisk Henseende, sagde, da han
havde læst Bogen, at mine politiske Meningsfæller vist vilde blive
skuffede ved den. Denne Ytring er mere berettiget end Dr.
Brandes' forvrængende Spot. Ganske vist, overfor en Anskuelse
som »aristokratisk Radikalisme« (der, i Forbigaaende sagt, hellere
kunde kalde sig »radikal Aristokratisme«) vilde selv den største
% --- p. 148 ---
\setcounter{footnote}{0}%
\opage{148}Grad af Upartiskhed, jeg vilde kunne hæve mig til, ikke være
nok til at afføde en anerkendende Dom.

Lige saa usandt som \emph{det} er, naar det siges om mig, at
»jeg godkender de bestaaende Institutioner«, lige saa lidet rammes
jeg ved den Bemærkning, at jeg føler mig tilfreds i min Tidsalder.
Tvært imod, jeg kunde ønske mig en langt lykkeligere Tid
at leve i. Tiden er baade for gammel og for ung, for overmoden
og for grøn. (De Fænomener, hvor det grønne og det
overmodne mødes i skøn Endrægtighed, er mig især vederstyggelige).
Der mangler Harmoni i Tro og i Liv. Overleveringen
tynger, og dog er det ikke formet og fuldbaarent, som skal
afløse den. Det gælder for Livet i den enkeltes Sind, og det
gælder for det sociale Liv. Men naar der (som Tilfældet er i en
Etik) er Tale om mulige Udgangspunkter for Viljen og Handlen,
ser jeg ikke, der er Grund til udelukkende at dvæle ved den Art
Betragtninger. Det gælder da at se til, hvad der kan faas ud
af det givne Grundlag. --- Enhver Tanke om Fremtiden og enhver
Indgriben er til syvende og sidst et Spring i Mørket, da vi ikke
kunne overse alle Forhold i deres Konsekvenser. Behageligere
er det maaske, og fornemmere (og der er jo dem, der »fremfor
alt elske det fornemme«), at stille sig udenfor Geled og lade
Tankerne svæve og Viljen hvile i det ubestemte eller i det fanta-
stiske. For dem, der (med Rette eller Urette) føle det som deres
Pligt at klare deres Tanker om de Opgaver, der ere os stillede
under de Forhold, som nu én Gang ere de, vort Liv er underlagt;
for dem maa den Forvisning være ledende, at Fremtiden
udvikler sig af, hvad der allerede nu er anlagt og til Stede som
Mulighed, og ikke af det, en lunefuld Fantasi kan foregøgle.
Fører hin Forvisning end ikke til at udtale glimrende og æstetisk
pointerede Tanker (eller rettere Ord), saa fører den maaske ad
sin jævne Vej til sundere og mere blivende Resultater.

2. Det er paa Tide at gaa \emph{over til} det andet af de to
Stridspunkter, som angaar Velfærdsmoralen.

Ligesom jeg i min første Artikel erklærede, at jeg ikke
bestemt kunde afgøre, i hvilken Betydning Sætningen om de
store Mennesker i Dr. Brandes' Artikel »Aristokratisk Radikalismes skulde tages, saaledes erklærede jeg fremdeles, at jeg ikke
ved de enkelte Udtalelser og Udviklinger bestemt kunde afgøre,
om det var Fr. Nietzsche eller G. Brandes eller dem begge, jeg
havde for mig. Dette faar særlig Anvendelse paa den \emph{Kritik af}
% --- p. 149 ---
\setcounter{footnote}{0}%
\opage{149}
\emph{Velfærdsmoralen, som} fandtes i Dr. Brandes' Artikel, og
som hørte til det i den, som overraskede mig mest. Jeg tog da
det Parti at holde mig til Indvendingerne mod Velfærdsmoralen,
som de forelaa, og omtalte deres Forfatter kun i ubestemte Udtryk
(f. Eks. »den, der har opstillet denne Indvending«). Men 
jeg kunde ganske vist ikke faa andet ud, end at Dr. Brandes,
om han end ikke selv havde fundet paa disse Indvendinger, saa
dog gav dem sin fulde Tilslutning, Derom syntes hele Fremstillingens
Art og Tone mig at vidne. Paa flere andre Punkter
tog desuden Dr. Brandes bestemt Afstand fra Nietzsche, i det
mindste ved smaa indskudte Vendinger, som det ogsaa her vilde
have været ham en let Sag at anbringe. Men her er (man efterse
»Tilskueren«s Augusthæfte, p. 584 f.) ikke en Stavelse, ikke en
Betoning, som antyder nogen Afstand mellem Referent og Forfatter.

Og jeg skal sige Dr. Brandes endnu en Grund, hvorfor jeg
antog, at han nu sluttede sig til Forkastelsen af Velfærdsmoralen,
skøndt jeg vidste, at han tidligere selv har hyldet denne. Jeg
\emph{forudsatte, at der dog var saa megen Konsekvens i
hans Tankegang}, at \emph{han ikke paa én Gang antog
Nietzsches Hovedpaastande og Grundanskuelser og
fastholdt Velfærdsmoralen}\footnote[1]{Dette var ogsaa en af Grundene til, at jeg saa skarpt kriticerede Dr. Brandes’ første Udtalelser om Stormennesker, skøndt deres Fortolkning var noget usikker. \emph{Ti den i Artiklen fremførte} og \emph{uimodsagte Forkastelse af Velfærdsmoralen forekom mig at tyde paa}, at Ordene \emph{om Stormenneskene som Historiens Maal maatte} tages i \emph{deres strængeste Betydning. Kun} ved den mildere Fortolkning, Dr. Brandes nu har givet, bliver det muligt at fastholde Velfærdsmoralen. Men saa maa man ogsaa erklære sig uenig med Nietzsche i alle Hovedpunkter. Dr. Brandes kan ikke slippe for den Modsigelse, jeg her finder hos ham, uden enten ved at overføre den paa Nietzsche selv (\emph{da den mildere} Fortolkning af Ordene om Stormenneskene er uforenelig med Velfærdsmoralens Forkastelse), eller ogsaa ved ganske at sige sig løs fra Nietzsches Lære (da den \emph{strængere} Fortolkning af Ordene om Stormenneskene er uforenelig med Antagelsen af Velfærdsmoralen).}. Fra Nietzsches Standpunkt
er det konsekvent at forkaste Velfærdsmoralen. Denne er jo kun
det filosofisk afklarede og formulerede Udtryk for en Tendens,
der (som jeg i min første Artikel p. 870 har omtalt) kan for%
% --- p. 150 ---
\setcounter{footnote}{0}%
\opage{150}følges tilbage til den gammelindiske Husordning, en Tendens,
der har udfoldet sig gennem græsk og kristelig Etik og ligger til
Grund i den moderne Humanitetsfølelse. Med hele denne Udvikling
bryder Nietzsche i sin »Moralgenealogi«, ja han bryder
ganske Staven over den. Og i Nutiden ser han den største Fare
i »de lidendes Sammensværgelse mod de sunde og sejerrige.«
Hans Feltraab kan udtrykkes saaledes: Ve eder, I lidende!
Fremad, I lykkelige! Derimod fastholder Velfærdsmoralen, hvad
jeg har kaldt de lidendes Menneskerettigheder.

Derfor vilde jeg ikke tro mine egne Øjne, da jeg i Dr.
Brandes' anden Artikel læste følgende: »Jeg havde sat mig til
Opgave at fremstille Nietzsche, ikke at markere \emph{ethvert nok
saa underordnet Punkt}\footnote[1]{Udhævet af mig.}, paa hvilket jeg var uenig med
ham. Jeg vilde derved ikke have opnaaet andet end at slaa
Skikkelsen i Stykker. Derfor gengav jeg hans Indvendinger mod
Lykkemoralen uden Kommentar.« Et underordnet Punkt! Det
er jo det, hvorpaa alt kommer an, saa sandt som det ved en
etisk Anskuelse især kommer an paa, hvilket Vurderingsprincip
man hylder. Skiller man sig fra Nietzsche paa dette Punkt, maa
man væsentlig ændre sin Opfattelse af ham paa flere andre
Punkter. Det stemmer meget vel med Velfærdsmoralen at tale
imod slap Filantropi og sygelig Medlidenhedsfølelse; men det er
uforeneligt med den at lovprise en Anskuelse som Nietzsches,
der hævder en Dualisme mellem de store og de smaa, de stærke
og de svage, de lykkelige og de ulykkelige, de sunde og de lidende.
Nietzsche har selv følt, at jo stærkere denne Modsætning
betones, des større bliver Konflikten med Velfærdsmoralen, og at
der maa træffes et Valg mellem Fornemhedens Standpunkt og
det Humanitetsstandpunkt, som Velfærdsmoralen betegner. Velfærdsprincipet
fører til \emph{en stadig fortsat} Stræben efter \emph{at udvide}
de lykkeliges Kreds; den Nietzscheske Grundtanke sætter en evig
Kløft mellem de lykkelige og de ulykkelige.

En ung Forfatter, Hr. Arnold Rohde, har da ogsaa truffet
sit Valg. Han har taget Dr. Brandes paa Ordet, proklameret
Nietzsche som Tidens Frelser, Brandes som hans Profet og sig
selv som en af de første Disciple. Han har med en vis Konsekvens
forkyndt den nye Lære under Titelen: »Aristokratisk
% --- p. 151 ---
\setcounter{footnote}{0}%
\opage{151}Radikalisme contra Velfærdsmoral«\footnote[1]{Siden jeg er kommen til at omtale Hr. Rohdes lille livlige Skrift, vil jeg i stor Korthed gaa ind paa hans Kritik af min Etik, der danner Indledning til hans begejstrede Tilslutning til den aristokratiske Radikalisme, hvilken han som sagt opfatter mere konsekvent end Dr. Brandes. Hans Kritik misforstaar min Fremstilling i tre Hovedpunkter. 1) Imod min Bedømmelse af, hvad jeg har kaldet Øjebliksstandpunktet, indvender han, at man jo dog i det enkelte Øjeblik kan tage Hensyn til hele Livet, sætte det enkelte Øjeblik over alt andet, hvad der kan opleves. Han har ikke lagt Mærke til, at saaledes som Øjebliksstandpunktet er beskrevet, og saaledes som det forsøgtes gennemført af Aristippos, udelukker det netop et saadant Hensyn (hvad enten dette fremtræder i Erindringens, Forventningens eller Instinktets Form). Det Standpunkt, han forsvarer, er altsaa et andet end det, jeg har omtalt. 2) Jeg lærer i min Etik, at \emph{Velfærd bestaar i Virksomhed, altsaa i Kultur} (hvorved jeg forstaar Bearbejdelsen af det naturligt givne). Hr. Rohde mener, at jeg derved opgiver Velfærdsprincipet. Og naar jeg saa føjer til, at \emph{al Kultur ikke uden videre er} af \emph{etisk} Værdi, fordi den kan komme i Strid med Velfærdsprincipet, saa mener han, at jeg „kører i Ring“. Han har glemt den elementære logiske Regel, at fordi A er B, behøver B ikke altid at være A.: En Tænker er et Menneske, men ikke alle Mennesker ere Tænkere. De etisk-sociale Opgaver, hele det sociale Spørgsmaal, opstaa netop ved, at Kulturudviklingen og den Arbejdsordning, den medfører, kommer i Strid med Velfærdsprincipets Fordringer. Kulturudviklingen kan føre til, at der lægges for store Byrder paa enkelte Klasser og Individer, at enkelte Ævner begunstiges paa andres Bekostning, at Kræfterne splittes paa usund Maade. De specielle etiske Problemer opstaa, som jeg i en udførlig Fremstilling af alle Sider af den materielle og den ideelle Kultur har søgt at vise, netop ved, at skøndt Velfærd bestaar i Kulturvirksomhed, fører ikke al Kulturvirksomhed (i den Betydning, i hvilken jeg tager Ordet) til; sand Velfærd. 3) Naar jeg taler om Slægtens Velfærd som etisk Maalestok, tillader Hr. Rohde sig uden videre at indskyde (se p. 13 i hans Pjece) Ordet „nulevende“, saa at han tillægger mig den Mening, at alt skal gaa ud paa at tilfredsstille \emph{den nulevende Slægts} Krav! --- Paa Grundlag af disse tre Misforstaaelser (foruden nogle mindre) er Hr. Rohde saa i en Avis bleven anerkendt som den, der har „paavist Principforvirringen“ i „Universitetets Patentmoral“! Ak nej, der doceres ogsaa en helt anden Moral ved Universitetet, og vi have nu tillige set, at ogsaa Dr. Brandes fastholder et Patent paa den af Hr. Rohde angrebne Lære.}%
. Men han bliver udtrykkelig
desavoueret af Dr. Brandes, overfor hvis ovenanførte Udtalelse
dette »contra« optræder med en ubønhørlig Ironi.

% --- p. 152 ---
\setcounter{footnote}{0}%
\opage{152}Naar jeg gør Regnskabet op med Dr. Brandes, faar jeg
da følgende Resultat.

\emph{Hvad det ene af de to Stridspunkter angaar (om
de store Mennesker som Historiens Maal), har han
brækket Spidsen af de dristige Sætninger, der gav
Artiklen »Aristokratisk Radikalisme« dens egentlige}
Betydning. Derved er dette Stridspunkt skaffet ud af Verden,
og der er i Artiklen, hvad det egentlige Tankeindhold angaar,
ikke sagt andet end hvad der ofte er sagt før. \emph{Hvad det
andet Punkt angaar (om Velfærdsmoralen), har han
hildet sig i en tydelig Selvmodsigelse og er havnet i
et Standpunkt, som ikke blot hans Modstandere, men
ogsaa hans Mester og hans Discipel erklære for uholdbart}.
--- Mon jeg ikke skulde have brugt et »mildt« Udtryk, naar
jeg har erklæret Artiklen om Nietzsche for ukritisk?

\asterism

Hermed \emph{kunde} jeg --- hvad Sagen angaar --- slutte Polemiken
med stor Sindsro. Men --- med Fare for, at Dr. Brandes
skal sige, at det ikke »klinger ægte«, kan jeg ikke lade være at
udtale, at lige saa modstræbende som jeg er gaaet til \emph{denne}
Strid, lige saa ilde til Mode har jeg følt mig ved at se en af de
betydeligste Mænd i vor Literatur optræde paa den Maade, det
her er sket. Jeg vil ikke undersøge, om Dr. Brandes er en stor
Mand, heller ikke, om det havde været bedre, om han havde
kunnet undværes. Men jeg vil slaa fast, at han ikke kunde
undværes her hjemme. Det skal aldrig glemmes ham, hvilket
Liv og hvilken Bevægelse han har vakt, og hvor mange friske
Strømninger han har ledet ind over vort Aandsliv. Han har
Egenskaber, som jeg vilde misunde ham, hvis jeg ikke i saa høj
Grad beundrede dem. (Hvem skulde vel mere kunne vurdere
Betydningen af Ævnen til at vække og begejstre, end den, hvis
Lærervirksomhed stadig bringer ham i Berøring med de unge,
og som saa ofte maa sige sig selv, at om han ellers skulde have
nogle Betingelser for sin Gerning, denne har han i hvert Tilfælde
ikke?). Vort Land maa være meget rigt paa store Kræfter,
siden det har haft Raad til ikke at tage denne betydelige Kraft i
% --- p. 153 ---
\setcounter{footnote}{0}%
\opage{153}sin Tjæneste. Der er gjort stor Uret mod Georg Brandes. Men
det forekommer mig rigtignok, at i de senere Aar er der ingen,
der øver større Uret mod ham end han selv. Maatte han dog
ikke fortære sin Kraft i usammenhængende Journalistik og i løse
Proklamationer af Fænomener, der især synes store, fordi de ere
barokke og sygelige! Maatte han kunne samle sig til et saadant
 betydeligt Værk, som vi endnu have Ret \emph{til at} vente det fra en
Haand som hans!
\signature{\textsc{Harald Høffding}.}

\end{document}
